\documentclass[a4paper]{article}
\usepackage[pdftex]{graphicx}
\usepackage[utf8]{inputenc}
\usepackage{enumerate}
\usepackage{icomma}
\usepackage{amssymb}
\usepackage{amsmath}
\usepackage{tikz}
\usepackage{href-ul}
\hypersetup{
	colorlinks=true,
	linkcolor=blue,
	urlcolor=blue}
\usepackage{geometry}
\geometry{a4paper, top=15mm, left=15mm, right=15mm, bottom=15mm,
	headsep=10mm, footskip=12mm}

% Start the document
\begin{document}

\begin{tikzpicture}(10,3)
	\draw[ultra thick](2,0) --(10,0) -- (10,1.3) --(2,1.3) -- (2,0);
	\draw[fill=black](2,0)-- (0,.6) -- (2,.6) -- (2,0);
	\node at (6,.6) {\large Lösungsblatt von \href{https://www.okuyakl.de}{www.okuyakl.de}};
\end{tikzpicture}
\vspace{0.5 cm}


\noindent{\bf Aufgabe 1.a)}\\
$$
\renewcommand{\arraystretch}{2}
\begin{array}{rcll}
\tan^2 x &=& {1 \over 4(1-\sin^2 x)} \\
\tan^2 x  &=& {1 \over 4\cos^2 x}  \\
 {\sin^2 x \over \cos^2 x} &=& {1 \over 4\cos^2 x}  &| \cdot \cos^2x\\
\sin^2 x &=& { 1 \over 4} &|\sqrt{\quad}\\
\sin x &=& \pm { 1 \over 2} \\
x_1 &=& {\pi \over 6}\\
x_2 &=& -{\pi \over 6}  \\
x_2 &=& -{\pi \over 6} + 2\pi &={11 \pi \over 6} \\
 x_3 &=&\pi-{\pi \over 6} &= {5\pi \over 6}\\
x_4 &=& \pi -\left(-{\pi \over 6}\right) &= {7 \pi \over 6}\\

\end{array}
$$

\noindent{\bf Aufgabe 1.b)}\\
$$
\renewcommand{\arraystretch}{2}
\begin{array}{rcl}
\sin x \cos x &=& \sqrt{1-\sin^2 x}\\
\sin x \cos x &=& \sqrt{\cos^2 x} \\
\sin x \cos x &=& \cos x \\
\sin x &=& 1 \quad \lor \quad \cos x =0\\
x_1 &=& {\pi \over 2} \\
x_2 &=& {3 \pi \over 2}
\end{array}
$$

\noindent{\bf Aufgabe 1.c)}\\
$$
\renewcommand{\arraystretch}{2}
\begin{array}{rcll}
{\cos^2 x \over \sin^2 x} &=& {3 \over 4 \sin^2 x}\\
\cos^2 x &=& {3 \over 4} &|\sqrt{\quad}\\
\cos x &=& \pm \sqrt{3 \over 4} \\
x_1 &=& {\pi \over 6} \\
x_2 &=& \pi-{\pi \over 6}   \\
x_2 &=& {5\pi \over 6} \\ 
2\pi - x_3 &=&  {\pi \over 6} \\
x_3 &=& { 11 \pi \over 6} \\
2\pi - x_4 &=&  {5\pi \over 6} \\
x_4 &=& { 7 \pi \over 6}
\end{array}
$$

\vspace{2 cm}

\noindent{\bf Aufgabe 1.d)}\\
$$
\renewcommand{\arraystretch}{2}
\begin{array}{rcll}
\sin x + {1 \over 2} &=& { 1 \over 2 \sin x} &|\cdot \sin x\\
\sin^2 x + {1 \over 2}\sin x  &=& { 1 \over 2} & | - {1 \over 2}\\
 \sin^2 x + {1 \over 2}\sin x  - { 1 \over 2} &=& 0 &| \sin x = u \\
u^2 + {1 \over 2} u - {1 \over 2} &=& 0 &\textnormal{Mitternachtsformel}\\
\Rightarrow u_1 &=& -1 \\
u_2 &=& {1 \over 2} \\
\sin x &=& -1 \\
x_1 &=& {3 \pi \over 2}\\
\sin x &=&  {1 \over 2} \\
x_2 &=& {\pi \over 6} \\
x_3 &=& {5 \pi \over 6}
\end{array}
$$

\noindent{\bf Aufgabe 1.e)}\\
$$
\renewcommand{\arraystretch}{1.5}
\begin{array}{rcll}
\sin^2 x + {1 \over 2}\sqrt{2}\sin x + \cos x &=&  \cos x &|-\cos x\\
\sin^2 x + {1 \over 2}\sqrt{2}\sin x  &=&  0\\
\sin x \left(\sin x +{1 \over 2}\sqrt{2}\right) &=&0\\
\sin x &=& 0 \\
x_1 &=& 0 \\
x_2 &=& \pi \\
x_3 &=&2 \pi \\
\sin x +{1 \over 2}\sqrt{2}&=&0\\
\sin x&=& - {1 \over 2}\sqrt{2} \\
x_4 &=& -{\pi \over 4} + 2 \pi &= {7 \pi \over 4}\\
x_5 &=& \pi - \left( -{\pi \over 4}\right) &= {5\pi \over 4}
\end{array}
$$

\noindent{\bf Aufgabe 2.}\\
a) $$-\sin x+ \sin x =0$$
b) $$2\cos x - \cos x =\cos x$$
c) $$\cos x + \cos x = 2 \cos x$$
d) $$ \sin x - \sin x =0$$
e)$$\cos x - \cos x=0$$
f) $$\sin x + \sin x = 2 \sin x$$
g)$$-\sin x + \sin x =0$$


\onecolumn
\noindent{\bf Aufgabe 3.}\\
a) $$ (\sin^2 x + \cos^2 x)^2 = 1^2 =1$$
b)$${{ \sin x \over \cos x}\over \sqrt{1+ { \sin^2 x \over \cos^2 x}}}={\sin x \over   \sqrt{\cos^2 x (1+ { \sin^2 x \over \cos^2 x})}}={\sin x \over \sqrt{\cos^2 x + \sin^2 x}}= {\sin x \over 1}=\sin x$$

\noindent{\bf Aufgabe 4.}\\
Weitere Ergebnisse im Intervall [$- \pi ; 5 \pi$] erhalten wir mit der Veranschaulichung der Sinusfunktion:

\includegraphics[width=12 cm]{sin144}

\noindent
Aufgrund der Symmetrie der Sinusfunktion gilt:
$$\sin{x}= \sin{(\pi - x)}$$ 
Also ist eine weitere Lösung der oberen Gleichung bei:
$$x_m= \pi + {\pi \over 5}= 1,2 \pi$$  
Für ganzzahlige $k \in \mathbb{Z}$ gilt aufgrund der Periodizität der Sinusfunktion:
$$ \sin{x}= \sin{(x+2k\pi)}$$
Also sind nun alle Lösungen
$$
\begin{array}{c|c|c|c|c}
k   & -1      & 0        & 1       & 2       \\
\hline
x_n & -       & -0,2 \pi & 1,8 \pi & 3,8 \pi \\
x_m & -0,8\pi &  1,2 \pi & 3,2 \pi &  -      \\
\end{array}
$$
\noindent {\bf Aufgabe 5.}\\

\begin{minipage}{0.5\textwidth}
	Wir formen zunächst den Ausdruck im Sinus um, damit die Verschiebung in x-Richtung ersichtlich 
	wird:
	$$ f(x)=2\sin{\left({1 \over 2} \left( x - {\pi \over 2}\right)\right)}$$
\end{minipage}
\begin{minipage}{0.5\textwidth}
	\includegraphics[width=7 cm]{f2sin144}
\end{minipage}

\begin{center}
	\includegraphics[width=7 cm]{../../viecher/endcomic.pdf}
	
	Hier geht es zurück zum \href{https://www.okuyakl.de/math/m10gozL025/aa025.pdf}{Aufgabenblatt}
\end{center}

\end{document}